% =============================================================================
\section{Overview}
\label{sec:overview}
% =============================================================================

\begin{figure}[t]
  \centering
  % Wider than the text block: the diagram bleeds 0.9cm into each margin (which
  % are 2.4cm) so the stage labels stay legible. Change the 18cm to \linewidth
  % to keep it strictly inside the text block.
  \makebox[\linewidth][c]{\resizebox{18cm}{!}{% =============================================================================
%  Pipeline overview figure -- TikZ source.
%
%  Included from sections/03-method-overview.tex via \input. Requires, in the
%  preamble of main.tex:
%
%      \usepackage{tikz}
%      \usetikzlibrary{arrows.meta,positioning,fit,backgrounds}
%
%  Coordinates are in the units set by x=/y= below; the whole figure is about
%  19.5 x 9.0 units, so at 0.78cm/unit it lands close to \textwidth on A4 with
%  2.4cm margins. Change ONE number (the x=/y= scale) to resize everything.
%
%  Layout was prototyped and visually checked before being written out here;
%  every coordinate below is the one that was verified.
% =============================================================================
\begin{tikzpicture}[
  x=0.78cm, y=0.78cm,
  font=\sffamily,
  % ---- boxes -------------------------------------------------------------
  stage/.style={
    draw=black!38, line width=0.5pt, fill=white, rounded corners=2.2pt,
    align=center, inner sep=2.6pt, font=\sffamily\scriptsize,
    minimum width=2.10cm, minimum height=0.84cm},
  pano/.style={stage, fill=black!4, minimum width=2.02cm, minimum height=0.74cm},
  anchorbox/.style={
    draw=black!72, line width=0.9pt, fill=black!7, rounded corners=2.6pt,
    align=center, inner sep=3pt, font=\sffamily\scriptsize\bfseries},
  outbox/.style={anchorbox, line width=1.2pt, fill=black!12},
  % ---- containers --------------------------------------------------------
  panogroup/.style={
    draw=black!38, line width=0.7pt, fill=black!2, rounded corners=3pt},
  lanelabel/.style={
    anchor=west, font=\sffamily\scriptsize\bfseries, text=black!62},
  grouplabel/.style={font=\sffamily\scriptsize\itshape, text=black!62},
  % ---- edges -------------------------------------------------------------
  flow/.style={draw=black!58, line width=0.5pt, -{Latex[length=3.4pt,width=2.6pt]}},
  sem/.style={flow, dashed, dash pattern=on 2.4pt off 1.8pt},
]

% ---- lane backgrounds (drawn first, so boxes sit on top) -------------------
% one \filldraw per lane: fill, border, y-span
\filldraw[fill=blue!8,   draw=blue!22,   line width=0.6pt, rounded corners=2.6pt] (5.55, 5.67) rectangle (13.05, 6.93);
\filldraw[fill=orange!10, draw=orange!30, line width=0.6pt, rounded corners=2.6pt] (5.55, 4.42) rectangle (13.05, 5.68);
\filldraw[fill=green!9,   draw=green!25,  line width=0.6pt, rounded corners=2.6pt] (5.55, 3.17) rectangle (13.05, 4.43);
\filldraw[fill=violet!8,  draw=violet!22, line width=0.6pt, rounded corners=2.6pt] (5.55, 1.92) rectangle (13.05, 3.18);
\filldraw[fill=yellow!14, draw=yellow!40, line width=0.6pt, rounded corners=2.6pt] (5.55, 0.67) rectangle (13.05, 1.93);
\filldraw[fill=black!5,   draw=black!18,  line width=0.6pt, rounded corners=2.6pt] (5.55,-0.68) rectangle (13.05, 0.58);

\node[lanelabel] at (5.71,6.74) {Terrain};
\node[lanelabel] at (5.71,5.49) {Objects};
\node[lanelabel] at (5.71,4.24) {Ground cover};
\node[lanelabel] at (5.71,2.99) {Water};
\node[lanelabel] at (5.71,1.74) {Sky};
\node[lanelabel] at (5.71,0.39) {Motion};

% ---- input ----------------------------------------------------------------
\node[anchorbox, minimum width=1.80cm, minimum height=1.05cm] (photo) at (0.80,3.10)
  {Single\\photograph};

% ---- three panorama variants ----------------------------------------------
\draw[panogroup] (2.08,1.72) rectangle (4.32,4.52);
\node[grouplabel] at (3.20,4.78) {Three panorama variants};
\node[pano] (porig) at (3.20,3.98) {Original\\panorama};
\node[pano] (pterr) at (3.20,3.06) {Object-removed\\panorama};
\node[pano] (psky)  at (3.20,2.14) {Sky-inpainted\\panorama};

\draw[flow] (1.70,3.10) -- (2.02,3.10);

% ---- lane contents ---------------------------------------------------------
% three-step lanes: x = 6.60, 9.30, 12.00 ; two-step lanes: x = 6.60, 12.00
\node[stage] (t1) at (6.60,6.25) {Height field};
\node[stage] (t2) at (9.30,6.25) {Constrained\\reconstruction};
\node[stage] (t3) at (12.00,6.25) {Mesh + texture};
\draw[flow] (t1) -- (t2);   \draw[flow] (t2) -- (t3);

\node[stage] (o1) at (6.60,5.00) {Detection\\+ typing};
\node[stage] (o2) at (9.30,5.00) {Appearance\\clustering};
\node[stage] (o3) at (12.00,5.00) {3D assets};
\draw[flow] (o1) -- (o2);   \draw[flow] (o2) -- (o3);

\node[stage] (g1) at (6.60,3.75) {Spatial\\statistics};
\node[stage] (g2) at (12.00,3.75) {Population\\synthesis};
\draw[flow] (g1) -- (g2);

\node[stage] (w1) at (6.60,2.50) {Region\\extraction};
\node[stage] (w2) at (12.00,2.50) {Levelled\\surface};
\draw[flow] (w1) -- (w2);

\node[stage] (s1) at (6.60,1.25) {Lighting\\estimation};
\node[stage] (s2) at (12.00,1.25) {Skybox};
\draw[flow] (s1) -- (s2);

\node[stage] (m1) at (6.60,-0.10) {Video\\generation};
\node[stage] (m2) at (9.30,-0.10) {Motion\\extraction};
\node[stage] (m3) at (12.00,-0.10) {Rigging};
\draw[flow] (m1) -- (m2);   \draw[flow] (m2) -- (m3);

% ---- panorama -> lanes (geometry) -----------------------------------------
% which variant feeds which lane is the load-bearing detail; see the report text
\draw[flow] (pterr.east) to[out=20,in=180]  (5.50,6.30);   % object-removed -> terrain
\draw[flow] (porig.east) to[out=-10,in=180] (5.50,5.05);   % original       -> objects
\draw[flow] (porig.east) to[out=-30,in=180] (5.50,3.80);   % original       -> ground cover
\draw[flow] (pterr.east) to[out=-25,in=180] (5.50,2.55);   % object-removed -> water
\draw[flow] (psky.east)  to[out=-20,in=180] (5.50,1.30);   % sky-inpainted  -> sky
\draw[flow] (porig.east) to[out=-55,in=180] (5.50,-0.05);  % original       -> motion

% ---- semantics (dashed) ----------------------------------------------------
\draw[sem] (porig.east) to[out=45,in=150] (t2.north west);   % region typing -> ridge anchoring
\draw[sem] (o2.south)   to[out=-90,in=70] (g1.north east);   % exemplars     -> population stats

% ---- convergence -----------------------------------------------------------
\node[anchorbox, minimum width=1.90cm, minimum height=1.10cm] (asm) at (14.55,3.10)
  {Scene\\assembly};
\foreach \yy in {6.25, 5.00, 3.75, 2.50, 1.25, -0.10}
  { \draw[flow] (13.10,\yy) to[out=0,in=180] (asm.west); }

\node[outbox, minimum width=2.40cm, minimum height=1.15cm] (out) at (17.60,3.10)
  {Explorable 3D\\environment};
\draw[flow] (asm) -- (out);

% ---- legend ----------------------------------------------------------------
\draw[draw=black!58, line width=0.5pt] (0.35,-1.35) -- (1.05,-1.35);
\node[anchor=west, font=\sffamily\scriptsize, text=black!62] at (1.15,-1.35)
  {geometry / assets};
\draw[draw=black!58, line width=0.5pt, dashed, dash pattern=on 2.4pt off 1.8pt]
  (3.55,-1.35) -- (4.25,-1.35);
\node[anchor=west, font=\sffamily\scriptsize, text=black!62] at (4.35,-1.35)
  {semantics};

\end{tikzpicture}
}}
  \caption{The pipeline decomposes a single image into five tracks with different
  physical roles. Solid arrows carry geometry, dashed arrows carry semantics.
  The three panorama variants at the left are maintained deliberately: geometry is
  derived from the object-removed panorama, semantics from the original.}
  \label{fig:pipeline}
\end{figure}

\subsection{Design principle: role-specialised decomposition}

The organising commitment of Into Frame is that a scene is not one thing. Ground,
discrete objects, ground cover, water and sky differ in how they must be
reconstructed, how they are represented, how they are textured, and how they
behave at runtime. Forcing them through one representation means accepting the
weakest common denominator.

\Cref{fig:pipeline} shows the resulting structure. Concretely, the pipeline
maintains five tracks:

\begin{description}[leftmargin=1.2em,style=nextline]
  \item[Terrain.] A single continuous height field, meshed with variable density,
  carrying collision. Reconstructed from calibrated panoramic depth and completed
  by a constrained solve.
  \item[Objects.] Discrete entities detected, classified, clustered by visual
  similarity, and reconstructed once per (class, appearance) group, then instanced
  at every detection's position.
  \item[Ground cover.] Populations too numerous to place individually as
  reconstructions, scattered by density over a mask derived from semantics and
  colour.
  \item[Water.] Regions typed as water, extracted as a separate levelled surface
  with an animated shader.
  \item[Sky.] The panorama above the horizon, inpainted where objects were removed,
  used as both skybox and image-based lighting source.
\end{description}

\subsection{Execution model}

The pipeline comprises 40 active stages. Each declares typed inputs and outputs
against a shared \emph{context}, and each may report whether its output already
exists. When it does, the stage is skipped. A stage that runs marks everything
downstream dirty, so a rerun is a linear cascade from the earliest changed stage.

This is a development-time contribution as much as an engineering one. Full
pipeline execution takes hours, so the ability to invalidate one stage and rerun
only its consequences is what makes iteration possible at all. It also creates a
specific hazard we encountered repeatedly: a stage's cache-validity predicate must
mirror the conditions under which the stage itself would skip work. Where the two
diverge, the stage reports incomplete forever and every downstream stage reruns on
every invocation.

A second consequence is that the configuration file is not part of the cache key.
Changing a parameter therefore does not invalidate the stage it belongs to, and an
explicit rerun directive is required. This is a deliberate trade --- hashing
configuration would invalidate far more than necessary --- but it is a sharp edge.

\subsection{Stage sequence}

\Cref{tab:stages} groups the stages by track; \cref{sec:implementation} reports
what each costs to run. The ordering is not arbitrary:
several dependencies are counterintuitive and were established by failure. Object
segmentation runs on the \emph{original} panorama, before foreground inpainting,
because inpainting removes the objects that must be detected; but terrain geometry
is derived from the \emph{inpainted} panorama, because occluders would otherwise
punch holes in the ground. Maintaining both variants and routing each consumer to
the correct one is a recurring structural theme.

\begin{table}[t]
\centering
\small
\caption{Pipeline stages grouped by reconstruction track.}
\label{tab:stages}
\begin{tabular}{@{}llp{7.4cm}@{}}
\toprule
\textbf{Track} & \textbf{Stages} & \textbf{Purpose} \\
\midrule
Panorama    & 1--4   & Caption, perspective depth, $360^\circ$ outpainting, super-resolution \\
Semantics   & 5--8   & Panoramic depth, object segmentation, classification, typing \\
Cleanup     & 9--11  & Foreground removal, equirectangular LoRA correction \\
Regions     & 12--15 & Semantic region typing (original and object-removed), lighting estimation, sky inpainting \\
Geometry    & 16--19 & Panoramic depth, metric calibration, height field, region map \\
Terrain     & 20--25 & Constrained reconstruction, erosion refinement, meshing, texturing \\
Objects     & 26--30 & Recognition, detection, instance refinement, correlation, clustering \\
Population  & 31--33 & Spatial statistics, population synthesis, ground cover \\
Assembly    & 34--35 & Asset generation, scene assembly and placement \\
Animation   & 38--42 & Video generation, motion extraction, rigging, animation \\
\bottomrule
\end{tabular}
\end{table}
